%%%%%%%%%%%%%%%%%%%%%%%%%%%%
\section{Binary Arithmetic}
%%%%%%%%%%%%%%%%%%%%%%%%%%%%
Binary arithmetic is the foundation of all digital computer systems because binary numbers, composed of only two digits (0 and 1), correspond directly to the on-off states of transistors, which are the basic building blocks of digital circuits. These two states are convenient for digital electronics because they can easily represent high and low voltage levels, aligning with binary logic \cite{knuth1997}. Despite its simplicity, binary arithmetic follows the same general rules as decimal arithmetic. This reduction to two possible states significantly simplifies hardware design and enables efficient processing \cite{rosen2018}.

%***************************
\subsection{Binary Addition}
%***************************
Binary addition operates similarly to decimal addition but involves only the digits 0 and 1, with carry rules adapted to binary logic. Binary addition requires only four basic cases:

\[
\begin{array}{l}
    0 + 0 = 0 \\
    0 + 1 = 1 \\
    1 + 0 = 1 \\
    1 + 1 = 10 \quad (\text{where 1 is carried to the next higher bit})
\end{array}
\]

When the sum of two bits equals 2 (in decimal), this triggers a carry to the next higher bit, as shown in the last case. This carry-over behavior is fundamental to binary addition and is analogous to the carry mechanism in decimal addition. For instance, adding the binary numbers \( 1011 \) and \( 1101 \) proceeds as follows:

\[
\begin{array}{cccccl}
  & 1& 1 & 1 & 1 &     (\text{Carry}) \\ % Carry row
  &  & 1 & 0 & 1 & 1    \\ % First binary number
+ &  & 1 & 1 & 0 & 1    \\ % Second binary number
  \hline
  & 1 & 1 & 0 & 0 & 0    \\ % Result row
\end{array} = 11000_2
\]

This example demonstrates that binary addition, although straightforward, can require multiple carry operations when the bit values sum to 1 in successive columns \cite{knuth1997art}.

%******************************
\subsection{Binary Subtraction}
%******************************
Binary subtraction is similar to decimal subtraction, with borrowing rules adapted to the binary system. Binary subtraction also involves four basic cases:

\[
\begin{array}{l}
    0 - 0 = 0 \\
    1 - 0 = 1 \\
    1 - 1 = 0 \\
    0 - 1 = 1 \quad (\text{borrow 1 from the next higher bit})
\end{array}
\]

When subtracting a larger bit from a smaller bit (such as \( 0 - 1 \)), we need to "borrow" from the next higher bit, just as in decimal subtraction. This process is slightly simplified in binary because the only possible values are 0 and 1, reducing the complexity of borrowing. For instance, consider subtracting \( 101 \) from \( 1100 \):

Step 1: Initial Setup. We add a 0 to the subtrahend to align the digits.
\[
\begin{array}{rcccc}
    & 1 & 1 & 0 & 0 \\ % First binary number (1100)
  - & 0 & 1 & 0 & 1 \\ % Second binary number (101)
  \hline
    &   &   &   &   \\
\end{array}
\]

Step 2: Subtract the rightmost column.

Here, \(0 - 1\) requires borrowing from the second column. Since the second column also has \(0\), we need to move leftward until we find a \(1\) to borrow. The third column has a 1. The minuend is set to \colorbox{yellow}{0} and the borrow row is set to two.
\[
\begin{array}{rcccc}
    \text{Borrow:} &   &   & 2 & 0 \\ % Borrow initiated for the rightmost column
                   & 1 & \fcolorbox{yellow}{yellow}{0} & 0 & 0 \\ % First binary number
                 - & 0 & 1 & 0 & 1 \\ % Second binary number
  \hline
                   &   &   &   & 1 \\ % Result after subtracting rightmost column
\end{array}
\]

Step 3. Move the \(2\) to the first column subtracting \(1\) from the second column. 
\[
\begin{array}{rcccc}
    \text{Borrow:} &   &   & 1 & 2 \\ % Borrow initiated for the rightmost column
                   & 1 & 0 & 0 & 0 \\ % First binary number
                 - & 0 & 1 & 0 & 1 \\ % Second binary number
  \hline
                   &   &   &   & 1 \\ % Result after subtracting rightmost column
\end{array}
\]



Step 4. We can subtract \(1-0=1\) and place the result in the second column.
\[
\begin{array}{rcccc}
    \text{Borrow:} &   &   & 1 & 2 \\ % Borrow initiated for the rightmost column
                   & 1 & 0 & 0 & 0 \\ % First binary number
                 - & 0 & 1 & 0 & 1 \\ % Second binary number
  \hline
                   &   &   & 1 & 1 \\ % Result after subtracting rightmost column
\end{array}
\]

Step 5. The third column has \(0-1\) so we need to borrow again. We set the fourth column minuend to \colorbox{yellow}{0} and place a 2 in the borrow row. We can then complete the subtraction \(2-1=1\) and place it in the third row of the result.
\[
\begin{array}{rcccc}
    \text{Borrow:} &   & 2 & 1 & 2 \\ % Borrow initiated for the rightmost column
                   & \fcolorbox{yellow}{yellow}{0} & 0 & 0 & 0 \\ % First binary number
                 - & 0 & 1 & 0 & 1 \\ % Second binary number
  \hline
                   &   & 1 & 1 & 1 \\ % Result after subtracting rightmost column
\end{array}
\]

Step 7. We complete the result by subtracting \(0-0=0\) and place a 0 in the fourth column of the result.
\[
\begin{array}{rcccc}
    \text{Borrow:} & 0 & 2 & 1 & 2 \\ % Borrow initiated for the rightmost column
                   & 0 & 0 & 0 & 0 \\ % First binary number
                 - & 0 & 1 & 0 & 1 \\ % Second binary number
  \hline
                   & 0 & 1 & 1 & 1 \\ % Result after subtracting rightmost column
\end{array}= 0111_2
\]

This example illustrates that borrowing in binary works similarly to decimal borrowing.

%***************************************
\subsection{Binary Multiplication - TBD}
%***************************************
\lipsum[1]

%*********************************
\subsection{Binary Division - TBD}
%*********************************
\lipsum[1]
